\documentclass{report}
\usepackage{amsmath,amssymb}
\setlength{\parindent}{0mm}
\setlength{\parskip}{1em}
\begin{document}
\begin{center}
\rule{6in}{1pt} \
{\large Joris Degroote \\
{\bf Solving time-dependent multi-physics problems using multiple instances of the same solvers}}

Ghent University \\ Department of Flow \\ Heat and Combustion Mechanics \\ Sint-Pietersnieuwstraat 41 \\ B-9000 Ghent \\ Belgium
\\
{\tt Joris.Degroote@UGent.be}\\
Robby Haelterman\\
Jan Vierendeels\end{center}

This research aims at accelerating the simulation of fluid-structure
interaction (FSI), which is a multi-physics problem involving the mutual
influence between a fluid flow and a deformable structure. In addition to
the governing equations in the fluid and structure domains, also two
equilibrium conditions have to be satisfied on the contact surface
between these domains, namely equality of velocity and traction on both
sides of the contact surface. Examples of fluid-structure interaction
problems are the motion of heart valves, the vibration of heat exchanger
tubes and the bending of wind turbine blades.

Fluid-structure interaction problems can be solved in a monolithic or a
partitioned way. The monolithic approach is to solve all equations
simultaneously, taking into account the interaction between the fluid and
the structure during the solution process. Conversely, the partitioned
approach is to solve the equations for the fluid and for the structure
separately and to perform coupling iterations until the equilibrium
conditions are satisfied. The main advantage of the partitioned approach
is that it allows to reuse existing solvers for the fluid flow and for
the structural deformation. In this research, the focus lies on
partitioned techniques and more specifically on quasi-Newton methods for
time-dependent simulations.

Traditionally, only one flow solver and one structural solver are used
for a fluid-structure interaction simulation. Each of them can run on
multiple cores in a parallel calculation in which the domain is split
over the different cores. However, there is typically a limit to the
number of cores which can be used to accelerate the simulation, depending
on the size of the problem, the computer hardware and the parallelization
of the software. Therefore, the main idea of this research is to use
multiple flow solvers and multiple structural solvers, all solving the
same nonlinear problem and each one of them possibly running on multiple
cores. Hence, an additional level of parallelization is obtained, besides
the domain decomposition.

The presented algorithm considers one of the flow solvers and one of the
structural solvers as the masters, which actually calculate the solution
of the nonlinear problem in the current time step using a quasi-Newton
algorithm. Obviously, the convergence of the quasi-Newton iterations is
faster if the approximation for the Jacobian is better. In the
quasi-Newton methods, this approximation is constructed using the
information obtained during the quasi-Newton steps, by considering
differences between consecutive iterations. The additional solvers
calculate the residual of the nonlinear problem in the current time step
after adding a step from the previous time step to the current value of
the solution. The resulting information in the current time step helps to
improve to approximation of the Jacobian and consequently accelerates the
coupling iterations between the master solvers. Clearly, this is only
useful for time-dependent problems.

As a test case, the propagation of a pressure wave through an
incompressible and inviscid fluid in a flexible tube is calculated. As
the number of flow solvers and structural solvers is increased from 1 to
8, the average number of coupling iterations per time step reduces from
approximately 8 to 3. Consequently, a speed up of the calculation with a
factor more than 2 is obtained, on top of the speed up due to
parallelization by domain decomposition. The results indicate a
relatively low parallel efficiency for this additional level of
parallelization, so increasing the number of solvers should only be
applied if it results in a higher speed up than increasing the number of
cores per solver. Furthermore, it has to be noted that the exact number
of coupling iterations depends on the system load as the additional
solvers can finish their calculation slightly before or after the master
solvers perform a quasi-Newton step, resulting in the inclusion of
additional information in the approximate Jacobian or not.

This so-called multi-solver quasi-Newton algorithm is only considered as
an initial step in the development of this kind of algorithms and
hopefully also other algorithms with multiple solvers will be developed.


\end{document}
